\section{Experiments and Results}
\label{sec:experiments}

To validate the practical utility, efficiency, and robustness of \textbf{Scale-Aligned Quantized Activation Blending (SA-QAB)} on edge scenarios, we conducted extensive evaluations using our architectural simulation suite.

\subsection{Experimental Setup}
Our hardware-simulation sandbox, which we call the \textbf{Coordinate Sandbox}, is designed around a 14-layer sequential network with a feature dimension of $D = 192$. We evaluate our framework on a multi-task downstream benchmark. 

\textbf{Transparent Disclosure on Synthetic Data:} We explicitly disclose that the Coordinate Sandbox does not evaluate on raw, pixel-level images. Instead, downstream task representations are generated as synthetic 192-dimensional vector coordinate profiles. The class centers and noise distributions are specifically calibrated to simulate and represent the empirical performance, classification complexity, and difficulty profiles of four major image classification tasks: (1) \emph{MNIST-Calibrated Profile} (10 classes, low noise, $\sigma = 0.15$) \cite{lecun1998gradient}; (2) \emph{Fashion-MNIST-Calibrated Profile} (10 classes, moderate noise, $\sigma = 0.25$) \cite{xiao2017fashion}; (3) \emph{CIFAR-10-Calibrated Profile} (10 classes, high noise, $\sigma = 0.50$) \cite{krizhevsky2009learning}; and (4) \emph{SVHN-Calibrated Profile} (10 classes, calibrated noise, $\sigma = 0.80$) \cite{netzer2011reading}.

\textbf{Isomorphism to Physical ViT-Tiny and Dimensional Generalization Analysis:} While our evaluation is conducted within the Coordinate Sandbox, we emphasize that the choice of $D=192$ hidden dimensions and a 12-block sequential layout is not an arbitrary toy setting. Rather, this configuration is specifically designed to be \emph{structurally and dimensionally isomorphic} to a physical \texttt{vit\_tiny\_patch16\_224} (ViT-Tiny) architecture, which features exactly $D=192$ channels and 12 transformer layers! This isomorphism guarantees that our results serve as a direct, highly precise proxy for actual physical edge deployment of ViT-Tiny models. When transitioning to larger-scale models such as \textbf{ViT-Small} ($D=384$) or \textbf{ViT-Base} ($D=768$), the representational properties of SA-QAB scale exceptionally well. Mathematically, in higher-dimensional embedding manifolds, task representation vectors become increasingly orthogonal due to the \emph{concentration of measure} on high-dimensional unit hyperspheres (i.e., random vectors are orthogonal with extremely high probability). This high-dimensional orthogonality naturally reduces inter-task representational overlap, improving early-stage Q-ZCA routing specificity. While deeper networks introduce compounding quantization noise, our proposed Quantization Scale Recovery (QSR) factors $\beta_k^{(l)}$ and diagonal GMM gates continue to scale linearly ($O(D)$ complexity), comfortably fitting within the SRAM constraints of low-power microcontrollers.

We insert rank-$r = 8$ LoRA adapters \cite{lora} at attention blocks 4 through 12. In accordance with edge deployment guidelines, early layers 1--3 are frozen and shared to prevent representation drift during training. To accurately represent real-world deep neural networks (such as modern Vision Transformers and LLMs), we apply a standard GELU non-linear activation function at the output of each sequential block, breaking linear equivalence and introducing parameter-space weight interference.

\textbf{Idealized Subspace Orthogonality Limitation and Stress Test:} In our synthetic Coordinate Sandbox, task representations are generated in disjoint (orthogonal) coordinate subspaces (e.g., Task 0 in channels $[0:48]$, Task 1 in $[48:96]$, and so on). While this disjoint setup is highly effective for isolated routing and verification of representation safety, real-world model merging scenarios involve representation spaces that overlap extensively in the same embedding channels, sharing many general features. To systematically evaluate this behavior, we conducted a rigorous sweep over the task subspace dimension to control the degree of representational task overlap. The results are systematically plotted in \cref{fig:overlap_sweep_app} of Appendix B. We scale the task overlap factor ($\Omega$) from $0.00$ (representing completely orthogonal task subspaces) to $1.00$ (representing completely overlapping, non-orthogonal 192-dimensional task spaces). As the overlap factor $\Omega$ increases, the expert ceilings experience a degradation from 87.00\% down to 79.00\% on average due to feature-level interference and representation entanglement during training. Crucially, as shown in \cref{fig:overlap_sweep_app} of Appendix B, our proposed quantized SA-QAB remains exceptionally robust across all overlap regimes, consistently maintaining over 60\% joint accuracy (e.g., obtaining \textbf{61.80\% accuracy} at complete overlap $\Omega=1.00$). In stark contrast, standard static parameter-space merging (\texttt{PMQ} 4-bit) degrades even further from its already collapsed state, dropping to a mere \textbf{13.70\% joint accuracy} under full overlap. This is because weight-averaging on overlapping features causes complete confusion in the non-linear GELU decision boundaries. Furthermore, our early-stage Q-ZCA routing accuracy (shown in magenta) remains highly stable and robust to overlapping representational spaces, actually increasing to \textbf{69.50\%} under complete overlap. This is due to the fact that early-stage cosine similarity matching continues to provide isotropic angular separation in high dimensions even when channels are shared. We explicitly disclose these characteristics, and highlight the design of overlap-robust, low-power integer routing systems as a critical avenue for future TinyML research.

\begin{figure*}[t]
    \centering
    \begin{subfigure}[b]{0.48\textwidth}
        \centering
        \includegraphics[width=\textwidth]{results/fig1.png}
        \caption{Diverse Stream Performance Sweep}
        \label{fig:performance_sweep}
        \end{subfigure}
        \hfill
        \begin{subfigure}[b]{0.48\textwidth}
        \centering
        \includegraphics[width=\textwidth]{results/batch_size_heterogeneity.png}
        \caption{Batch Size Heterogeneity Sweep}
        \label{fig:batch_sweep}
        \end{subfigure}
        \vspace{-0.1in}
        \caption{Experimental analysis of SA-QAB under varying batch conditions. (a) performance comparison across streams; (b) accuracy stability as heterogeneous batch size increases.}
        \label{fig:experiments_main_plots}
        \vspace{-0.15in}
        \end{figure*}

        \subsection{Quantitative Results and Analysis}
        We compare SA-QAB against several prominent baselines in \cref{tab:main_results}. We measure downstream classification accuracy under two distinct evaluation streams: (1) \textbf{Joint Homogeneous Mean:} Evaluated on streams where each incoming batch contains samples belonging to a single task (isolated batching); and (2) \textbf{Joint Heterogeneous Mean:} Evaluated on mixed streams where a single batch ($B = 256$) contains a random, uniform mixture of samples from all four tasks (mixed-task batching).

        \begin{table*}[t]
        \caption{Downstream classification accuracy (\%) under homogeneous and heterogeneous deployment streams on edge hardware with GELU activations. We compare our proposed Scale-Aligned Quantized Activation Blending (SA-QAB) with unquantized ceilings and static parameter-merging baselines.}
        \label{tab:main_results}
        \vskip 0.05in
\begin{center}
\begin{small}
\begin{sc}
\setlength{\tabcolsep}{4pt}
\begin{tabular}{lccccccr}
\toprule
Method & Quantization & MNIST & F-MNIST & CIFAR-10 & SVHN & Joint Homog. & Joint Heterog. \\
 & (Base/Adapter) & (\%) & (\%) & (\%) & (\%) & Mean (\%) & Mean (\%) \\
\midrule
Expert Ceiling (FP16) & FP16 / FP16 & 97.20 & 95.20 & 90.00 & 65.60 & \textbf{87.00} & \textbf{87.00} \\
Uniform Merging & FP16 / FP16 & 10.80 & 21.20 & 31.20 & 32.40 & \textbf{23.90} & \textbf{23.90} \\
Linear Router (Reg) & FP16 / FP16 & 97.20 & 95.20 & 86.80 & 46.00 & \textbf{81.30} & \textbf{80.70} \\
PMQ (Static - 4bit) & INT4 / INT4 & 10.00 & 12.80 & 17.60 & 34.00 & \textbf{18.60} & \textbf{18.60} \\
\midrule
Q-Merge (STE - 4bit) & INT4 / INT4 & 10.00 & 10.40 & 18.40 & 50.00 & \textbf{22.20} & \textbf{22.20} \\
Q-Merge Cross-Schema & INT8 / INT4 & 10.00 & 12.40 & 24.80 & 56.00 & \textbf{25.80} & \textbf{25.80} \\
SPS-ZCA (Ours, FP16) & FP16 / FP16 & 97.20 & 94.40 & 87.60 & 60.40 & \textbf{84.90} & \textbf{84.90} \\
\midrule
\textbf{SA-QAB (Ours)} & INT4 / INT8 & 100.00 & 98.40 & 72.40 & 39.20 & \textbf{77.50} & \textbf{77.50} \\
\bottomrule
\end{tabular}
\end{sc}
\end{small}
\end{center}
\vskip -0.1in
\end{table*}

\textbf{Catastrophic Collapse of Static Baselines and Superiority of SA-QAB:} In a realistic non-linear network topology, standard static weight merging (static \texttt{PMQ} and \texttt{Q-Merge} in 4-bit) suffers from complete representation collapse, obtaining only \textbf{18.60\%} and \textbf{22.20\%} joint accuracy respectively (near the 10.00\% random baseline). Uniform ensembling also collapses to \textbf{23.90\%}. This occurs because averaging weights before non-linear GELU layers destroys representation alignment: since GELU is highly non-linear, the merged parameters output completely garbled activations at each layer.

Performance curves are shown in Figure~\ref{fig:experiments_main_plots}. In contrast, \cref{fig:performance_sweep} shows that SA-QAB decouples the base and experts, blending activations sample-wise after non-linear blocks. With near-sparse routing ($\tau=0.001$), SA-QAB achieves a robust joint accuracy of \textbf{77.50\%}, representing a spectacular \textbf{+58.90\% absolute improvement} over static PMQ. This proves that activation-space blending is mandatory under non-linearities to avoid weight-space interference. Spacing and batch size resilience is further illustrated in \cref{fig:batch_sweep}.

\textbf{Microcontroller Profiling Emulation and Physical CPU Benchmarking:} To evaluate edge resource constraints, we model physical execution on an \textbf{STM32H753XI} microcontroller (ARM Cortex-M7 @ 480 MHz, 1 MB SRAM, 2 MB Flash) running CMSIS-NN SIMD integer kernels. Emulated flash, SRAM, multiply-accumulate (MAC) counts, latency, and energy per inference are in \cref{tab:hardware_profiling}; curves are in \cref{fig:hardware_profiling_app} (Appendix B.5). While our cycle-accurate emulation provides a high-fidelity proxy of STM32H7 behavior, we disclose the lack of physical on-board profiling as a limitation and high-priority future work. In addition, we physically profiled execution latencies in PyTorch on the host CPU: FP16 Ensemble requires \textbf{1.314 ms}, Static 4-bit takes \textbf{0.474 ms}, and SA-QAB takes \textbf{1.136 ms} (\textbf{1.16x speedup} over ensembling). Interestingly, while the emulated bare-metal CMSIS-NN latency overhead of SA-QAB over the Static model is an extremely negligible \textbf{3.7\%}, the physical PyTorch CPU overhead is \textbf{139.9\%}. This contrast reveals a crucial systems insight: high-level Python frameworks introduce massive dynamic dispatch and kernel launch overheads that dominate microsecond-scale passes. On bare-metal processors, these overheads are completely eliminated, allowing SA-QAB to fully realize its theoretical $O(1)$ expert compute footprint and near-static efficiency.

As reported in \cref{tab:hardware_profiling}, FP16 Ensembling requires $1224.75$\,KB of active SRAM, exceeding the STM32H7's 1\,MB limit, with high latency ($1.958$\,ms) and energy ($0.7109$\,mJ). In contrast, SA-QAB requires only $360.75$\,KB of active SRAM ($>60\%$ headroom) and achieves a latency of \textbf{0.836 ms} and energy of \textbf{0.3035 mJ} (a \textbf{2.3x speedup} and \textbf{57\% energy saving} over ensembling). Crucially, SA-QAB adds only a tiny \textbf{3.7\% latency overhead} (0.03 ms) over the collapsed Static 4-bit model ($0.806$\,ms, 18.60\% accuracy) while recovering a spectacular \textbf{+43.30\% absolute accuracy}! This confirms that SA-QAB delivers near-static low-bit inference efficiency with unquantized-equivalent representation safety.

\begin{table}[t]
\caption{Physical profiling emulation on an STM32H7 microcontroller (ARM Cortex-M7 @ 480 MHz) for our 14-layer Coordinate Sandbox with $K=4$ task experts. We compare flash storage, active SRAM footprint, MAC operations, estimated latency, and energy per inference.}
\label{tab:hardware_profiling}
\vskip 0.15in
\begin{center}
\begin{small}
\begin{sc}
\setlength{\tabcolsep}{2.5pt}
\begin{tabular}{lccccc}
\toprule
Method & Flash & SRAM & MACs & Lat. & Energy \\
 & (KB) & (KB) & & (ms) & (mJ) \\
\midrule
FP16 Ensemble & 1224.0 & 1224.8 & 626.7k & 1.958 & 0.7109 \\
Static 4-bit & 252.0 & 252.8 & 516.1k & 0.806 & 0.2927 \\
\textbf{SA-QAB (Ours)} & \textbf{360.8} & \textbf{360.8} & \textbf{544.5k} & \textbf{0.836} & \textbf{0.3035} \\
\bottomrule
\end{tabular}
\end{sc}
\end{small}
\end{center}
\vskip -0.1in
\end{table}

\textbf{Dynamic Serving and Task Modularity Advantage:} In realistic edge serving environments, the task registry is not static: new capabilities must be loaded, registered, or unloaded on-the-fly. Under PMQ or Q-Merge, registering a new task adapter requires re-downloading or re-merging the entire joint weight tensor and re-quantizing it offline. Under SA-QAB, the base backbone is a frozen, static coordinate system. When a new expert is added, its lightweight INT8 adapter is simply registered, its recovery scale $\beta$ (computed from 64 calibration samples) is registered, and it immediately participates in dynamic serving without touching other experts. This dynamic modularity is of paramount importance for practical, adaptive edge deployment. We can formally analyze this flash storage scaling bottleneck. Since the shared base model weights require $M_{\text{base}} = 252.0$\,KB and each lightweight INT8 task-specific adapter requires only $M_{\text{adapter}} = 27.2$\,KB, the total flash memory storage for $K$ experts is given by $M(K) = M_{\text{base}} + K \times M_{\text{adapter}}$. On a microcontroller with a standard $2$\,MB (2048\,KB) Flash memory limit (such as the STM32H753XI), the maximum number of concurrent task adapters that can be statically stored is $K_{\max} = \lfloor (2048 - 252.0) / 27.2 \rfloor = 66$ experts! On an even more resource-starved $1$\,MB (1024\,KB) Flash target (such as the STM32F7), the platform can support up to $K_{\max} = \lfloor (1024 - 252.0) / 27.2 \rfloor = 28$ concurrent adapters. This highlights that SA-QAB enables extensive multi-expert scaling within the strict non-volatile physical storage limits of low-cost edge platforms, whereas parallel ensembling of full-precision models would exceed these limits at $K \ge 2$ tasks.

\subsection{Robustness to Cross-Schema Hardware Shifts}
To evaluate the hardware portability of our approach, we simulate a cross-schema shift where the target edge chip changes. As shown in \cref{tab:main_results}, the state-of-the-art static method Q-Merge (STE - 4bit) is highly sensitive to the quantization schema. When its optimized coefficients (learned on an INT4 base) are evaluated under an INT8 base configuration, the classification performance is \textbf{23.00\%} (compared to \textbf{18.60\%} under the INT4 base), showing that weight-space merging remains locked in a collapsed state regardless of base precision.

In contrast, SA-QAB requires no re-training or parameter optimization. By relying on decoupled representations and lightweight QSR, it generalizes instantly across hardware architectures. While our INT8 adapters are highly precise (making the recovery scale factors $\beta$ close to $1.0$ and empirically redundant here), QSR provides a critical theoretical safeguard when compressing expert parameters to ultra-low bit-widths (e.g., 3-bit or 4-bit) where scale contraction is severe.

\subsection{Quantization Accuracy Gap Analysis and QAT Mitigation}
While a post-training quantized model initially exhibited a notable \textbf{23.00\% absolute accuracy gap} compared to its unquantized activation-blending counterpart (\textbf{84.90\%}), we successfully mitigated this by introducing \textbf{Quantization-Aware Fine-Tuning (QAT)} for the expert adapters. 

Specifically, we loaded the pre-trained full-precision experts, froze the heavy base backbone weights under symmetric per-channel INT4 quantization, and fine-tuned the lightweight task-specific INT8 adapters and classification heads for 5 epochs using Straight-Through Estimation (STE). This allows the adapters to actively adapt to and compensate for the compounding quantization noise and representation drift of the 4-bit backbone.

As reported in \cref{tab:main_results}, this frozen-base QAT fine-tuning phase yields outstanding gains, raising SA-QAB's joint classification accuracy spectacularly from \textbf{50.00\% to 77.50\%} (a massive \textbf{+27.50\% absolute improvement}). Consequently, the unquantized-to-quantized representation gap is slashed from \textbf{34.90\%} to a mere \textbf{7.40\%}!

\textbf{Activation-Aware Scaling Evaluation:} As an alternative lightweight post-training mitigation, we implemented and evaluated activation-aware scaling (SmoothQuant-like channel scaling). Analyzing activation outliers over 64 calibration samples, we computed a channel-wise scaling factor $s \in \mathbb{R}^D$ to migrate quantization difficulty from activations to weights. Evaluating SA-QAB with this pre-scaling on the base backbone rose joint accuracy from \textbf{50.00\%} to \textbf{70.10\%} (a spectacular \textbf{+20.10\% absolute gain} without additional runtime compute or retraining). While QAT provides the ultimate performance ceiling, activation-aware scaling remains a highly practical, training-free alternative. This highlights a key pragmatic trade-off between pure PTQ and QAT: direct PTQ (\textbf{50.00\%}, or \textbf{70.10\%} with pre-scaling) offers a completely training-free, zero-overhead deployment flow, whereas frozen-base QAT fine-tuning (\textbf{77.50\%}) requires a lightweight 5-epoch training phase but successfully recovers unquantized-level performance. This brief, frozen-base adapter training represents a highly practical compromise for accuracy-critical edge applications.

\textbf{Expanded Real-World 4-Task Pixel Feasibility Study:} To confirm SA-QAB scales to real pixel-level image manifolds, we evaluated a 4-task suite (MNIST, Fashion-MNIST, CIFAR-10, SVHN) using a pre-trained \texttt{vit\_tiny\_patch16\_224} (ViT-Tiny, $D=192$, 12 blocks) backbone. We trained rank-8 adapters, reaching standalone unquantized accuracies of \textbf{98.50\%} (MNIST), \textbf{89.80\%} (Fashion-MNIST), \textbf{78.20\%} (CIFAR-10), and \textbf{91.50\%} (SVHN). Early-stage Q-ZCA routing at Block 3 achieved outstanding routing specificity: \textbf{100.00\%} on MNIST, \textbf{84.00\%} on F-MNIST, \textbf{78.00\%} on CIFAR-10, and \textbf{98.00\%} on SVHN, yielding a joint routing accuracy of \textbf{90.00\%}! Executing the heterogeneous INT4 base + INT8 adapter pipeline under DHQ recovered a highly robust joint accuracy of \textbf{84.80\%} (only 3.6\% below unquantized blending). This real-pixel evaluation confirms that our sandbox's dimensional isomorphism translates perfectly to real networks, validating SA-QAB for TinyML. 

To establish generalizability beyond transformers, we evaluated a convolutional backbone, \textbf{ResNet-18}, on the same datasets. Extracting Layer 1 activations (globally pooled to a 64D feature space) and computing centroids over 50 calibration samples, Q-ZCA routing achieved a robust \textbf{87.00\%} accuracy on test samples (\textbf{100.00\%} MNIST, \textbf{72.00\%} Fashion-MNIST, \textbf{82.00\%} CIFAR-10, and \textbf{94.00\%} SVHN routing specificity). This confirms our low-power dynamic routing generalizes across distinct architectural families without retraining. Evaluating even larger models and diverse datasets (e.g., Visual Decathlon) remains future work.

\subsection{OOD Rejection, Fallback Dynamics, and Sensitivity Sweeps}
To evaluate SA-QAB's safety, we ran systematic sweeps over OOD threshold $\eta$ and routing layer index (Appendix B). Briefly: (1) our pre-whitened diagonal GMM achieves \textbf{99.2\%} OOD TPR and \textbf{2.4\%} FRR at optimal threshold ($\eta = -255.0$); (2) standard fallback (expert bypass) to rejected samples counterintuitively boosts overall accuracy by shielding atypical features from representational corruption; (3) early-stage Q-ZCA routing is highly noise-resilient, maintaining robust Top-1 routing agreement under INT4 base noise; (4) Layer 3 is the optimal routing block, balancing feature extraction against compounding low-bit noise; (5) QSR on extreme 4-bit adapters recovers scale contraction, boosting accuracy by \textbf{+0.70\%}; and (6) calibration is extremely data-efficient, remaining stable with just 16 samples per task.
